% !TeX root = ../main.tex
% Chapter 9 -- Conclusion and Future Work
% Self-contained section file: no preamble, no \documentclass.
% HONESTY RULE: no experiments have been run. This is a proposal / methodology
% document. This chapter restates the (planned) contribution, reaffirms the
% honesty boundary, and lays out future work / re-add triggers. It contains NO
% measured results of our own. Any number it mentions belongs to the prior
% literature and is attributed as such.
% Mirrors the template's Chapter 5: \S5.1 conclusion + \S5.4 recommendations.

\chapter{Conclusion and Future Work}
\label{ch:conclusion}

\begin{quote}\small\itshape
No experiments have been run. This chapter draws no empirical conclusions about
the performance of our instrument; it has none to draw. It restates what the
project \emph{proposes} to deliver, reaffirms the boundary between a proposal and
a validated study, and pre-registers the conditions under which descoped work
would be re-added. Every quantitative value it mentions belongs to the cited prior
literature and is attributed to its authors as a background anchor.
\end{quote}

This chapter mirrors the template's concluding chapter: a restatement of the
contribution (\S\ref{sec:concl-contribution}, in the spirit of the template's
\S5.1 conclusion), a reaffirmation of the honesty boundary
(\S\ref{sec:concl-honesty}), and a future-work and re-add-trigger schedule
(\S\ref{sec:concl-future}, in the spirit of the template's \S5.4 recommendations
for further study). It is deliberately short: a proposal earns its conclusion not
by summarising findings but by stating precisely what it would take to convert the
proposal into a validated claim.

%====================================================================
\section{Restatement of the Contribution}
\label{sec:concl-contribution}

This work set out under a single discipline: \emph{not to build another cat-pain
detector}. The danger was never that such a detector is hard --- the literature
already reports strong numbers, including the validated Feline Grimace Scale (FGS)
anchor of Evangelista et al.\ \cite{evangelista2019} and the landmark-to-geometry
automated pipeline of Steagall et al.\ \cite{steagall2023} --- but that a frozen
backbone with a binary pain head and a handful of standard reliability metrics
would be \emph{indistinguishable} from prior work once a reviewer mentally swapped
the backbone and added an extra metric. We therefore relocated the contribution
away from the engine and onto two transportable methods and one welfare-aware
operating discipline, and we conceded the engine in writing.

The thesis the document defends is stated in
Chapter~\ref{ch:introduction} and recapitulated here in its load-bearing form:
this is not another cat-pain detector because its headline deliverables are
\textbf{two transportable methods that no prior feline-pain work produced} ---

\begin{enumerate}
  \item a \textbf{per-AU \mbox{VLM-as-AU-rater} agreement protocol}: a protocol for
  measuring whether a frozen vision--language model can weak-label FGS action units at
  human-rater agreement (the result is pending an independent per-AU veterinary anchor,
  and measures weak-labeling capability only if that anchor is independent and the VLM
  rubric differs from the vet's --- otherwise it measures rubric-following), reported as
  five per-AU quadratic $\kappa$ values
  \emph{gated on the confidence-interval lower bound} against a veterinary anchor,
  with released code that runs the same comparison on any face corpus; and
  \item a \textbf{capture-condition confound-attribution protocol}
  (FGS-BG-Gap plus per-AU energy-based pointing-game faithfulness), framed
  one-directionally --- well-powered to \emph{detect} a capture-condition confound,
  explicitly under-powered to \emph{rule one out} --- and reusable on the next
  dataset rather than a verdict about this one ---
\end{enumerate}

\noindent shipped on an \textbf{explicitly-conceded \mbox{DINOv2$+$CORN} engine}
\cite{oquab2023dinov2,shi2023corn} that we never claim as novel, and wrapped in an
\textbf{honest binary-plus-abstention floor}: a welfare-asymmetric operating point
selected at a fixed pain-recall $\geq 0.90$ (anchored to
Evangelista et al.\ \cite{evangelista2019}) with the undertreat-to-overtreat harm
ratio swept as a range over a decision curve, and a defer-to-vet boundary reported
as a one-sided $95\%$ negative-predictive-value lower bound. The instrument is
framed throughout as decision-support triage --- ``grimace consistent with pain,
$X/10$; recommend veterinary assessment'' --- never as an autonomous analgesia
trigger.

Two consequences of this framing are themselves part of the contribution. First,
the engine concession is permanent and deliberate: a conceded
swap-the-backbone delta is harmless, whereas an oversold one is the reviewer's
favourite kill. Second --- and this is the strongest honesty commitment in the
project --- the \textbf{graded $0$--$10$ FGS layer ships as an
inspected-not-validated artifact}. The CORN ordinal heads are trained, decoded
distributionally, and inspected qualitatively, but the word ``graded'' is struck
from every validated-claim sentence in the document, and the abstract is required
to hold with the graded layer dropped entirely. The thing we cannot validate at
the available sample size, we do not assert. What survives that subtraction --- a
calibrated binary decision-support instrument, two portable methods, and a
welfare-aware operating discipline --- is the contribution, and it is not another
cat-pain detector.

%====================================================================
\section{Reaffirmation of the Honesty Boundary}
\label{sec:concl-honesty}

The single most important sentence in this document is that \textbf{no experiments
have been run}. This is a research proposal and a methodology, not a completed
empirical study. The document contains no measured accuracies, no $\kappa$ values
of our own, no calibration numbers, no abstention curves with our data on them,
and no figures of our own curves. Where a number appears --- the FGS validation
figures of Evangelista et al.\ (AUC~$0.94$, sensitivity~$90.7\%$,
specificity~$86.6\%$, threshold $>0.39$) \cite{evangelista2019}, the
$95.5\%$ landmark-geometry classification accuracy of Steagall et al.\
\cite{steagall2023}, or the landmark normalised-mean-error range of the CatFLW
benchmark \cite{martvel2024catflw} --- it belongs to the cited prior literature
and is presented strictly as a background anchor, never as a result of this work.

The boundary is enforced structurally rather than rhetorically. Chapter~\ref{ch:expected-outcomes}
replaces a results chapter with an \emph{expected-outcomes and gate-conditioned}
analysis: for each gate $G_0$ through $G_6$ (plus the weak-label reliability pilot
$G_{1\text{-}B}$) the document pre-registers the decision rule and the
\textsc{go}/\textsc{pivot}/\textsc{collapse} branch \emph{before} any datum is
seen, so that the eventual study cannot move its own goalposts. Three further
firewalls make the proposal falsifiable in advance: sensitivity and specificity at
the $0.39$ decision boundary are to be estimated \emph{only} on vet-confirmed
labels (the $\kappa$-against-VLM number is never treated as validation); the
cat-disjoint test split is frozen and hashed so that no individual straddles a
fold and the distinct-pain-cat denominator is always printed with
Clopper--Pearson or bootstrap intervals; and the word ``guaranteed'' is banned
from the abstention curve, which is reported only as a finite-sample lower bound.

Read correctly, then, the conclusions of this document are conditional and
methodological. They are claims about \emph{what would be measured, by what rule,
and what would ship in each branch} --- not claims about what was measured. The
abstract holds with the graded layer dropped; the contribution holds even if the
weak-label reliability pilot forces the planned pivot to the binary spine; and no
sentence in the document asserts a performance number that the project has not yet
earned the right to assert.

%====================================================================
\section{Future Work and Re-Add Triggers}
\label{sec:concl-future}

A proposal that descopes aggressively owes the reader an explicit account of
\emph{what was cut and what would bring it back}. The project's critical path was
narrowed to what is solo-deliverable on a single Apple-silicon machine with one
veterinary annotator; several capabilities were removed not because they are
undesirable but because the available sample size cannot support them honestly. We
therefore record each descoped item as a future-work line paired with a
pre-registered re-add trigger, so that a successor (a v2 effort, a larger team, or
a follow-on with a richer data agreement) inherits a schedule rather than a wish.
Table~\ref{tab:concl-readd} states the triggers; the prose below names the four
most consequential.

\begin{table}[t]
\centering
\caption[Descoped capabilities and their pre-registered re-add triggers]{Descoped
capabilities and the pre-registered conditions under which each would re-enter the
critical path. None is on the solo, single-vet timeline of v1; each is a v2
candidate gated on a concrete, measurable condition rather than on optimism. ``AU''
denotes a Feline Grimace Scale action unit; ``NPV'' denotes negative predictive
value.}
\label{tab:concl-readd}
\small
\begin{tabular}{@{}p{0.42\linewidth} p{0.52\linewidth}@{}}
\toprule
\textbf{Descoped capability} & \textbf{Re-add trigger (pre-registered)} \\
\midrule
Validated graded $0$--$10$ FGS claim
  & An \emph{independent} per-AU $0/1/2$ gold held-out set arrives \emph{and} the
    severe ($\mathrm{AU}{=}2$) cells reach double-digit counts. \\
\addlinespace
External request-only FGS dataset as a graded anchor
  & A request-only set clears a Data Use Agreement with a bounded turnaround;
    treated as a v2 windfall, never a v1 blocker. \\
\addlinespace
``Guaranteed'' NPV language on the abstention curve
  & A power calculation shows the budgeted veterinary $n$ certifies
    $\mathrm{NPV} \geq 0.90$ at an abstention rate $\leq 40\%$ (a larger vet
    account). \\
\addlinespace
Empirical $5\times 5$ cross-AU error-correlation matrix
  & The vet-labelled $n$ reaches the low hundreds with non-single-digit cells;
    until then a two-point $\rho$ sensitivity band stands in. \\
\addlinespace
Per-cell severe ($\mathrm{AU}{=}2$) calibration number
  & The anchor yields double-digit $\mathrm{AU}{=}2$ cells; until then the high
    end of the scale is collapsed. \\
\addlinespace
Engine novelty claim (\mbox{DINOv2$+$CORN})
  & Never --- conceded as plumbing permanently. \\
\bottomrule
\end{tabular}
\end{table}

\paragraph{Validating the graded layer.}
The graded $0$--$10$ regression is the natural next step, and converting it from
an inspected artifact into a validated claim is the project's principal future
work --- a move that mirrors the template's own recommendation to extend a binary
classifier toward a validated graded scale. Two conditions must both hold: an
\emph{independent} per-AU $0/1/2$ gold held-out set must exist (CatFLW
\cite{martvel2024catflw} cannot serve this role, being a landmark and
bounding-box resource rather than a graded-severity anchor), and the severe
$\mathrm{AU}{=}2$ cells must reach double-digit counts so that per-cell sensitivity
is estimable rather than spanning the whole unit interval. Only when both are met
does the project move from a binary spine to a validated graded regression with
the welfare-aware operating point carried forward onto the ordinal output.

\paragraph{A bounded data agreement as a v2 windfall.}
The richest accelerant would be a request-only feline-FGS dataset clearing a Data
Use Agreement with a bounded turnaround. Because such agreements carry unbounded
schedule risk on a solo timeline, the v1 plan deliberately keeps them off the
critical path; in v2 a cleared agreement is a windfall that simultaneously
supplies an external cohort for the confound-attribution protocol and a candidate
graded anchor. It is a re-add trigger, never a blocker.

\paragraph{Certifying the abstention floor.}
The defer-to-vet boundary currently ships as a one-sided $95\%$ NPV lower-bound
curve precisely because the budgeted veterinary sample cannot certify a useful
$\mathrm{NPV}$ floor at a tolerable abstention rate; certifying
$\mathrm{NPV} \geq 0.95$ would collapse the abstention rate toward zero and yield a
vacuous instrument. A larger veterinary account --- enough to certify
$\mathrm{NPV} \geq 0.90$ at an abstention rate of at most $40\%$ --- is the
re-add trigger that would let the abstention curve be reported as a certified floor
rather than an exploratory lower bound, and the word ``guaranteed'' remains banned
until that power calculation succeeds.

\paragraph{An empirical cross-AU correlation matrix.}
The propagation of per-AU error into the $0$--$10$ sum is currently handled by a
two-point $\rho$ sensitivity band (a \textsc{go} is declared only if the decision
holds under both $\rho{=}0$ and a pinned $\rho{=}0.3$), because a noisily fit
$5\times 5$ correlation matrix at the available $n$ launders sampling noise as
data-driven rigor and need not even be positive-definite. Once the vet-labelled
sample reaches the low hundreds with non-single-digit cells, the pinned band is
replaced by an empirical cross-AU error-correlation matrix --- a strictly more
honest input that requires only more data, not a new method.

\bigskip
\noindent In sum, the modal product this proposal plans to ship --- a calibrated
binary decision-support instrument with a welfare-asymmetric decision curve and a
one-sided-$95\%$-NPV abstention curve, two portable methods, and an
inspected-not-validated graded layer on a conceded engine --- is solo-deliverable
and stands on its own. The future work above is not a list of things missing from
that product but a schedule of well-defined, sample-size-gated upgrades, each of
which converts an honestly-deferred claim into a validated one the moment its
pre-registered trigger fires.
